\documentclass[11pt]{scrartcl}
\usepackage[a4paper,margin=2.4cm]{geometry}
\usepackage{../common/rlstudy}

\title{Week 14 --- Extended Reading: Offline RL}
\author{STUDY-RL}
\date{\today}

\begin{document}
\maketitle

This companion reads alongside \texttt{notes.pdf}. Each entry: full citation,
one-paragraph summary, ``Why it matters for this week'', and an annotated key
passage.

\section{Levine, Kumar, Tucker, Fu (2020) --- Offline RL Tutorial}
\textsc{S.\ Levine, A.\ Kumar, G.\ Tucker, J.\ Fu.} \emph{Offline Reinforcement
Learning: Tutorial, Review, and Perspectives on Open Problems.} arXiv:2005.01643, 2020.

\noindent\textbf{Summary.} The canonical reference for the field. Formalises the
problem; surveys distribution shift, importance sampling, BCQ/BEAR/CQL/AWR;
proposes a taxonomy of methods by where they place the constraint (policy vs.\
critic vs.\ density). Long but dense; chapters 4 and 6 are mandatory.

\noindent\textbf{Why it matters for this week.} It is the canonical mental model:
all subsequent papers are positioned with respect to this taxonomy.

\begin{quote}
``The fundamental challenge in offline RL is one of distributional shift: while we
optimise the policy under one distribution (the dataset $\mathcal{D}$), the policy
will visit a different distribution at evaluation time.''
\end{quote}
\emph{Commentary.} This one sentence is the entire field. Every algorithm in the
literature is a different answer to: \emph{which divergence between $d^\pi$ and
$d^{\pi_\beta}$ are we bounding, and how?}

\section{Kumar et al.\ (2020) --- Conservative Q-Learning}
\textsc{A.\ Kumar, A.\ Zhou, G.\ Tucker, S.\ Levine.} \emph{Conservative Q-Learning
for Offline RL.} NeurIPS 2020, arXiv:2006.04779.

\noindent\textbf{Summary.} Introduce the now-standard ``learned $Q$ should be a
lower bound on $V^\pi$'' criterion. The CQL penalty
$\log \sum_a \exp Q(s,a) - \E_{\pi_\beta} Q(s,a)$ added to Bellman regression
yields a lower-bound estimator under mild assumptions on the data coverage.
Strong on D4RL-MuJoCo.

\noindent\textbf{Why it matters.} It moved the field from ``policy constraint''
to ``value constraint''. Cleaner objective, fewer moving parts than BCQ/BEAR.

\begin{quote}
``By learning a conservative value function, CQL prevents the policy from
exploiting overestimation errors during the policy improvement step.''
\end{quote}
\emph{Commentary.} ``Prevents'' is overselling: it dampens. Tuning $\alpha$ is
where production bites; in practice we use the Lagrangian variant.

\section{Kostrikov et al.\ (2021) --- Implicit Q-Learning}
\textsc{I.\ Kostrikov, A.\ Nair, S.\ Levine.} \emph{Offline RL with Implicit
Q-Learning.} arXiv:2110.06169, 2021.

\noindent\textbf{Summary.} Replace explicit Q-on-out-of-support sampling with an
\emph{expectile} regression on the value: $V$ regresses to the upper expectile
of $Q(s, a \sim \pi_\beta)$, so it estimates ``best in-support'' without ever
naming the out-of-support action. Pairs with AWR/AWAC for policy extraction.
Faster, simpler, competitive with CQL.

\noindent\textbf{Why it matters.} IQL is the current default in many real-world
deployments precisely because there are no OOD samples needed at training time
--- the action space can be arbitrary.

\begin{quote}
``Our key insight is that we can perform the policy improvement step implicitly
by treating the value function as the upper expectile of the action-value
function under the data distribution.''
\end{quote}
\emph{Commentary.} The word ``implicitly'' carries the whole idea: the policy
is never used to query Q at unseen actions. Compare to AWAC, which still samples
actions from $\pi$.

\section{Tarasov et al.\ (2023) --- ReBRAC}
\textsc{D.\ Tarasov, V.\ Kurenkov, A.\ Nikulin, S.\ Kolesnikov.} \emph{Revisiting
the Minimalist Approach to Offline RL.} arXiv:2305.09836.

\noindent\textbf{Summary.} A careful re-engineering of TD3+BC: higher discount,
layer-norm in critics, per-step BC weight tuned to $\E[Q]$. No new ideas, just
correct execution. Sets the strongest non-transformer scores on D4RL as of late 2023.

\noindent\textbf{Why it matters.} Lesson: many offline RL ``algorithmic'' gains were
hyperparameter gains. Useful sanity check before reaching for the heavy tools.

\begin{quote}
``Many of the recent advances in offline RL can be attributed to a few well-chosen
implementation details rather than to the algorithmic novelty.''
\end{quote}
\emph{Commentary.} A polite rebuke of the field. Cite this paper whenever a new
SOTA claim arrives without ablations.

\section{Voloshin et al.\ (2021) --- OPE Empirical Study}
\textsc{C.\ Voloshin, H.\ Le, N.\ Jiang, Y.\ Yue.} \emph{Empirical Study of Off-Policy
Policy Evaluation.} NeurIPS Datasets and Benchmarks, 2021, arXiv:1911.06854.

\noindent\textbf{Summary.} A systematic comparison of IS, weighted-IS, doubly robust,
FQE, MAGIC, DICE estimators across a portfolio of MDPs. Verdict: \textbf{FQE} is the
single most reliable estimator when a function approximator is available; SNIPS is
the best ``no model'' baseline; doubly-robust requires careful nuisance tuning.

\noindent\textbf{Why it matters.} OPE is what stands between an offline-trained
policy and a production deployment. This paper is the ``which estimator should I use''
cheat-sheet.

\begin{quote}
``No single OPE method dominates across all settings; the choice depends on
horizon length, state-space dimensionality, and how well the behaviour policy is
known.''
\end{quote}
\emph{Commentary.} The horizon-length collapse is the key practical point: above
roughly 50 steps, IS-family estimators are useless; FQE or DICE is mandatory.

\vfill
\noindent\emph{Sourced from} \texttt{research/offline\_modelbased\_sota.md}
(track A, sections A.1--A.5, A.15--A.16).

\end{document}
